\section{Consolidation of the current project sources}
\label{sec:sr-source-review}
This section records the source review behind the present consolidation.
Its purpose is to distinguish a proved internal deduction, an imported
formal theorem and an assumption still awaiting a prose proof.
The review is of the files available on 21 November 2026.
Version identifiers and verified byte hashes are recorded in the
accompanying \srcid{short-route-sources.json}.

\subsection{The paper consolidation and historical review}
The
\href{https://principia-math.com/projects/d16ab288ed353bc5b10e292ace0ddbc1/paper}{project papers page}
has one current RH working manuscript from this account.
The earlier entry-point paper and the two earlier versions of the
translation were consolidated into this manuscript and their separate
attachments removed. Their source snapshots and metadata are preserved in
\href{https://github.com/ChristianKudeba/alpoge-rh-verification/blob/main/translation/archive/principia-before-194-page-consolidation.json}{the historical archive}.
Our separate Masullo audit remains a separate paper because it addresses
a different proof attempt. Other contributors' papers were left unchanged.
Before this consolidation the following five PDFs were identified and
reviewed; the older artifact links below are historical references.
\begin{description}
\item[Formal entry points.]
\href{https://principia-math.com/api/ws/files/402cfc03e76bc4f9eafd0eb930a80a20}{\srcid{alpoge-rh-formal-entry-points.pdf}}
records endpoint declarations and their formal dependencies.
Its six pages are not a six-page mathematical proof.
\item[Earlier working translation.]
\href{https://principia-math.com/api/ws/files/259cf87468cc745b1e81e358086a251e}{\srcid{alpoge-rh-working-translation.pdf}}
is an earlier version of our exposition, rather than an independent
argument that can discharge its successor's assumptions.
\item[The 163-page working translation.]
\href{https://principia-math.com/api/ws/files/8ca901b46a918a9293164d6a63d1ea7e}{\srcid{alpoge-rh-working-translation-163pp.pdf}}
is the supporting text incorporated in Part II.
Proposition~\ref{prop:sr-PB} now connects its printed minus descent to the
short route; no proof term is needed to make that connection.
\item[The separate analytic audit.]
\href{https://principia-math.com/api/ws/files/b561b940f8dd3d3b8f8a2d1d83b359dc}{\srcid{masullo-zero-count-audit.pdf}}
contains counterexamples and checks on the second proposed approach.
Section~\ref{sec:sr-shared-distribution} includes the checks relevant here.
\item[Masullo's draft.]
\href{https://principia-math.com/api/ws/files/e8e3c2bb97f9bf324d74bff487bfb435}{\emph{The Music of the Primes}}
is the 940-page version already examined in that audit.
Its prime-side normalization is comparable with
\eqref{eq:sr-prime-distribution}; its global spectral identification
cannot be used here without replacing the affected steps.
\end{description}
This review found no additional standalone proof closing the stated
inputs among these five files. The bytes of every listed PDF were hashed for this review.
In particular the Masullo version has not silently been replaced by
a revision resolving the identified issues.

\subsection{Anton's conductor-one reductions}
The first new specialization examined here was in the collaborative editor:
\href{https://principia-math.com/projects/d16ab288ed353bc5b10e292ace0ddbc1/editor?file=b64c07f1e23480429eacd6b354f04f26}{\srcid{ConductorOneRHCombined.lean}}.
The earlier 552-line snapshot has SHA-256
\[
 \begin{gathered}
 \texttt{e2bab4de455929689f2977514ae8a84b8}\\
 \texttt{004fbbbb19619cf7b6866d77c93d33e}.
 \end{gathered}
\]
The two displayed strings concatenate to one 64-digit hash.
For that snapshot, the author reports a successful Lean 4.33.1
compilation in 664 seconds.
The present consolidation independently checked the source bytes and
read the theorem signatures; it did not independently repeat that
compilation. The reported 33 direct imports reach 3,857 project
modules, besides Mathlib. The preceding 1,873-line version reached
3,856. Thus the shorter source is a useful specialization and
reorganization; its imported mathematical development remains extensive.

The top theorem is
\srcid{Principia.Ordinary.ConductorOneRH.rh}.
It is an implication with three explicit inputs.
The first is the oddness of the class number of every cyclotomic field
of \(2\)-power conductor; this is proved in
Lemma~\ref{lem:sr-oneprime} and its application above.
The other two inputs are as follows.

\paragraph{A uniform plus-length bound.}
In the source this is the predicate \srcid{Height}.
It requires a single natural exponent \(C\) and two fixed auxiliary
primes, before the host and the twist range are chosen.
For every sufficiently large principal host \(h\), every
\(t\geq2\) with \(h.\mathrm{Big}(t)\), and its CM field \(E\), it gives
\begin{equation}\label{eq:sr-principal-height}
 \sum_{\substack{\ell\in S\\\ell\ \mathrm{odd}}}\log\ell
       \sum_{a=1}^{t-1}s^+_{\ell,E}(a)
                       \leq(Nt)^C
       \quad\hbox{for every finite prime set }S .
\end{equation}
The host is for the trivial character. Its field and \(N\) are those
in Section~\ref{sec:sr-arithmetic}.
The size condition includes \(2t+1<q,d\) and excludes finitely many
fixed character-dependent constants. The two auxiliary-prime
inequalities and these fixed exclusions must hold for the same host.
Equation~\eqref{eq:sr-principal-height} is not a consequence of the minus
bound, since it controls the complementary conjugation eigenspaces.
Allowing \(C\) to depend on the individual host would change the assertion.

\paragraph{A construction of connected families.}
The final theorem also assumes the implication
\begin{quote}\raggedright
\srcid{HostExW (3 / 2) (GChi PChar.one).req}
\(\Longrightarrow\)
\srcid{FamilyClassConnected ((GChi PChar.one).cls (3 / 2))}.
\end{quote}
These are the archive's host-window and connected-family predicates.
They include geometric witnesses and their uniform estimates.
A prose deduction using this route must define those witnesses and prove
their properties. Naming the predicates does not do that.
Definition~\ref{def:sr-supply} instead states the exact signed-data
existence assertion used by the fully printed comparison in Part I.
The two interfaces should not be identified without constructing the
map between their witnesses.

The 552-line snapshot proves the following assembly within the imported
development. Its degree-one specialization removes the exceptional
odd-prime term. The minus estimate, the supplied dyadic estimate and
\eqref{eq:sr-principal-height} give the full strict-length budget.
The high-twist size theorem and the rank bound provide an exceptional
set with at most 24 elements. Finite-exception avoidance supplies a
common interval. This produces a principal cohomological window.
The assumed connected-family construction then gives the prime-sum
inequality, which excludes right-half-plane zeros. Modulus-one
Dirichlet \(L\) is \(\zeta\), so reflection gives ordinary RH.
The degree-one cancellation and the final identification of \(L\)
remove bookkeeping required for nontrivial characters.

Both combined revisions use
\[
 J_{\mathrm{archive}}=2^{26}(2C_KC_t/\epsilon)^{24}.
\]
The direct interval proof in \eqref{eq:sr-J} gives
\(J_{\mathrm{direct}}=2^{25}(2C_KC_t/\epsilon)^{24}\) under its
stated exceptional-set hypotheses.
These are different interval estimates. The sharper numerical
calculation does not itself establish those hypotheses.

\paragraph{Subsequent additions.}
The
\href{https://github.com/antoshashakov/riemann-hypothesis/blob/d85b449fa08d46836c934a7809fe3c9cae2e4a6a/verification/LogCubeRH.lean}{630-line cubic-logarithmic revision}
has the same direct import list and adds a weaker height option:
\eqref{eq:sr-logcube-height}.
Its pinned GitHub bytes match the updated Principia editor body,
with SHA-256
\[
 \begin{gathered}
 \texttt{d183191cac31584b7fa051847a87fbce3}\\
 \texttt{dac7ec4a44a9f5701fe4ad2f9904321}.
 \end{gathered}
\]
The author's report states that Lean checked both RH endpoints and
ten axiom reports. We checked the source bytes and its implications,
without repeating that Lean build.
Section~\ref{sec:sr-logcube} proves the new numerical absorption.

The further
\href{https://github.com/antoshashakov/riemann-hypothesis/blob/442bb98b2aa223ac4f28ec86e754e1a54b3ae081/verification/PrincipalHeightCounts.lean}{principal height counts},
\href{https://github.com/antoshashakov/riemann-hypothesis/blob/b8afa4ee2aca3cb3478b1f706cac6b03b258b9fe/verification/RegularCharacterHeight.lean}{regular-character vanishing}
and
\href{https://github.com/antoshashakov/riemann-hypothesis/blob/76ffabc62ab67b126afc128013ab6efa6f7d2a7c/verification/TwoCharacterCounts.lean}{two-character reduction}
have 342, 282 and 356 lines respectively.
Their uniform generic sum and nontrivial pointwise estimate remain
explicit inputs.
Section~\ref{sec:sr-two-counts} writes these inputs in mathematical
notation, proves the needed trivial-character vanishing and derives
the height bound from them with all prime weights retained.
This removes a separate CM-host assumption at 3 and 7.
It does not remove the need to prove the two rational estimates.

\subsection{New certificate constructions}
The newer repository work at
\href{https://github.com/ChristianKudeba/alpoge-rh-verification/tree/6186b8d5c272c08c90f6f5cce1a25d22232ed78d/verification/research/certificate-construction}{the pinned certificate-construction snapshot}
supplies explicit integer witnesses before the field realizations and
coefficient primes are chosen.
Its odd-minus certificates, Weber proof and dyadic-window adapters
assemble proved archive inputs. Its nonresonant plus certificate
uses \(\mathbb Q(\sqrt{-3})\), whose elementary arithmetic packet
is proved in Section~\ref{sec:sr-three-auxiliary}.
A pair of distinct auxiliary primes supplies two nonzero integers;
their squared product covers the one excluded prime of either
individual estimate. Products over a fixed set of rational character
orbits then give one integer independent of the later prime set.

The snapshot's notes explicitly retain the substantial
\srcid{PropGK} integer producer and the nontrivial resonant estimate
as inputs. They also identify \(\ell=3,a=1\) as a case that a
rank-four estimate stated only for \(a\geq2\) cannot cover.
Thus these constructions improve the arithmetic assembly without
providing the missing standalone counting proof.
The recorded kernel replay reports acceptance of the closed
odd-minus certificate, Weber theorem and the implication from a joint
window to the final cohomological window, using only the three standard
axioms. We inspected those records and the relevant sources; we did not
repeat their replay in this document build.
The joint-window producer still has mathematical content to translate.

\subsection{What imported results still need an internal proof}
The short file supplies the high-twist theorem by assembling the
class-counting ingredients. Its source lines 182--189 invoke
\srcid{lemmaK_of_steps} together with chain, prism, intersection,
theta, class, torsion and quaternionic bounds.
The explicit Poitou--Tate and local duality parameters are therefore
not its only mathematical content: the other ingredients are
imported proved results. The standalone manuscript still needs their
applicable statements and proofs, or a valid replacement for them.

The connected-family branch we inspected is
\href{https://github.com/ChristianKudeba/alpoge-rh-verification/blob/alpoge-source-7cbf7a0c/src/Development/AutomorphicArithmetic/ThetaCorrespondence/ThetaPairLifts/ThirdPlaceSeries/TopShape.lean}{\srcid{ThirdPlaceSeries.TopShape}}.
Its concluding comparison sets \(e_D=2,e_L=0\), turning the
exponent \(1/2+(e_D+e_L+7)\epsilon\) into \(1/2+9\epsilon\).
That substitution is elementary.
The hypotheses supplying the integral theta realizations, rank and
determinant bounds, genus mass, metric control, admissible twists and
connected geometric model are not elementary consequences of that
substitution. The longer route's supporting chapters establish several
comparisons from specified local data, while retaining the construction
of those data as inputs.

The community's generated
\href{https://principia-math.com/projects/d16ab288ed353bc5b10e292ace0ddbc1/editor?file=fd6fc09854184b61bca8dd800a4859b5}{\srcid{rh-unpacked.tex}}
earlier snapshot reports 1,067 translated modules and 1,001 reviewed modules
out of 175,187. We inspected relevant geometric and cohomological
entries in that snapshot. Subsequent project notes at 3:15 AM report
1,155 translated modules and 1,071 reviewed modules.
For example
\srcid{CuspidalModerateGrowthRapidDecay} separates the height used in
derivative growth from the quotient height used in its conclusion.
Its proof uses a cusp-scale decomposition, a unipotent-variation bound
and an iterated-gain theorem.
The following \srcid{CuspFormsRapidDecay} entry is a wrapper converting
that finite-derivative estimate into rapid decay.
These entries help track the proof, but appending their statements
would not prove the imported cusp-scale results or the uniform metric
estimates needed for the present families.
Likewise a published-hypothesis entry states an input; it is not an
independent proof of the input.

\subsection{The resulting proof ledger}
The odd prime-to-degree minus bound is now an internal consequence:
the strict characterwise descent, its derivative system and local
identities, the common Bernoulli integer and
Proposition~\ref{prop:sr-PB} form one chain.
The dyadic class-number base, the uniform 24-generator bound of
Proposition~\ref{prop:sr-dyadic-generators}, the elementary field degree,
the finite interval selection and numerical absorption are also proved here.
Section~\ref{sec:sr-dyadic-bounds} now proves a uniform free-lift
congruence, cokernel growth from depth seven, the matrix-logarithm
center bound, the 24-element exceptional set fixed on a whole horizon,
and the fixed-depth discriminant estimate. It also computes all local
cohomology corrections exactly and proves the resulting growth bound
of 25 per layer. Section~\ref{sec:sr-ideal-lattices} proves the
\srcid{SStein} fractional-ideal comparison and its covolume identity.
These are internal ingredients of the class-counting chain; its
remaining global estimates still require proofs.
The dyadic common-window
theorem is an internal deduction from \(K_2(C_K)\), not a proof of
that remaining high-twist input.
The analytic endpoint and signed finite comparison have complete
printed proofs. The new rational vanishing arguments and the assembly
from the two character estimates are printed in
Section~\ref{sec:sr-two-counts}.

To obtain an unconditional standalone proof by this route, the
following work remains.
\begin{enumerate}
\item Prove the uniform plus estimate
\eqref{eq:sr-principal-height}, or the weaker bound
\eqref{eq:sr-logcube-height}, with one constant before all hosts and
twist ranges. The two estimates in
Definition~\ref{def:sr-character-inputs} suffice for the polynomial
version, by Proposition~\ref{prop:sr-two-counts}. The height/deformation and automorphic counting
comparisons in Part II must be supplied with the construction and
compatibility inputs they currently assume.
\item Prove \(K_2(C_K)\) of Definition~\ref{def:sr-K2} with one
constant uniform in the principal pair and the twist.
The growth, center, exceptional-set and discriminant arguments of
Section~\ref{sec:sr-dyadic-bounds} then give the full dyadic
common-window bound by Theorem~\ref{thm:sr-dyadic-window}.
The remaining content is the quantitative high-twist
class-counting theorem, rather than the construction of the exceptional
set or its interval selection.
\item Construct the connected geometric families, or directly the
signed data of Definition~\ref{def:sr-supply}, on the required common
windows. This includes the relevant global theta identity, ramified
matching, cusp/metric and integral-lattice estimates with compatible
normalizations and constants.
\end{enumerate}
